\chapter{Introduction}
\label{chap:intro}
\chaptermark{Optional running chapter heading}
\section{Background \& Relevance}
\label{sec:background_and_relevance}

Index out of bounds errors are a well known reason for critical bugs appearing
in modern software, mainly occuring in places where the checks for proper indexing 
are either insufficient or lacking entirely. The most notable example of these errors
is access to an element of array that doesn't exist at the index provided.
\\
Since a lot of programming languages rely on the array as a commonly used data structure,
it is easy to encounter such an error in any part of the developed software.
The examples include but are not limited to color conversions [need citation \href{https://github.com/flutter/flutter/issues/89584}{89584} \href{https://github.com/abdulhadinaeem/widget-recoding-plugin/blob/1.0.4/CHANGELOG.md}{changelog}],
pixel values [need citation \href{https://github.com/fluttercommunity/flutter-draggable-scrollbar/issues/8}{8}] or handling of empty lists.
\\
Unfortunately, due to the nature of how array accesses happen, these errors are quite difficult to notice until they reach 
runtime or sometimes even production environments completely unnoticed.


\section{Problem Statement}
\label{sec:problem_statement}
Dart programming language is prone to index out of bounds errors. The current type system, 
though incressing robustness [citation dart language specs], does not currently support  
collection-length constraints. This makes developers rely on assumptions and runtime testing.
\\
This paper attempts to resolve this issue by introducing a refinement type system to statically
confirm safety of array access operations while adding little annotation burden to the developers.

\section{Objectives}
\label{sec:objectives}
We aim to achieve the following goals:
\begin{itemize}
    \item An implementation of an analyzer for a custom predicate language 
    \item An implementation of a refinement type checker and type inference algorithm with usage of Z3 constraints solver
    \item Integration of said checker into Dart analysis workflow \cite{noauthor_analyzer_nodate}
    \item Evaluation of the effectivness of the implementation on a package corpus from \href{pub.dev}{pub.dev}
\end{itemize}

\section{Limitations \& Scope}
\label{sec:limitations_and_scope}
The implementation will cover the following features of Dart and will allow users to write the following refinements:
\begin{itemize}
    \item Refinement type definitions that will allow users to extend integers, booleans and lists with refinements
    \item Definition of variables with refinement types
    \item Definition of recursive/non-recursive functions
    \item Usage of $length(xs)$ function in refinements for List type and its refinement extensions
\end{itemize}

% \section{Spacing \& Type}
% \label{sec:section}

% This is a section. This is a citation without brackets. and this is one with brackets \cite{A}. Multiple \cite{A,B,C} Here's a reference to a subsection: \ref{sec:subsection}. Citation of an online article \cite{D}. Citation of an online proceeding \cite{F}. The body of the text and abstract must be double-spaced except for footnotes or long quotations. Fonts such as Times Roman, Bookman, New Century Schoolbook, Garamond, Palatine, and Courier are acceptable and commonly found on most computers. The same type must be used throughout the body of the text. The font size must be 10 point or larger and footnotes\footnote{This is a footnote.} must be two sizes smaller than the text\footnote{This is another footnote.} but no smaller than eight points. Chapter, section, or other headings should be of a consistent font and size throughout the ETD, as should labels for illustrations, charts, and figures.

% \subsection{Creating a Subsection}
% \label{sec:subsection}

% \subsubsection{Creating a Subsubsection}
% \subsubsection{Creating a Subsubsection}
% \subsubsection{Creating a Subsubsection}

% \paragraph{This is a heading level below subsubsection}

% And this is a quote: 
% %
% \begin{quote}
% \blindtext
% \end{quote}

% \begin{figure}[hbt]
% \centering
% \includegraphics[]{figs/images.png}
% \caption{One kernel at $x_s$ (\emph{dotted kernel}) or two kernels at
% $x_i$ and $x_j$ (\textit{left and right}) lead to the same summed estimate
% at $x_s$. This shows a figure consisting of different types of
% lines. Elements of the figure described in the caption should be set in
% italics, in parentheses, as shown in this sample caption.}
% \label{fig:example}
% \end{figure}

% This is a table:
% % currsize is not set in the long table environment, so we need to set it before we set it up.
% \makeatletter
% \let\@currsize\normalsize
% \makeatother

% % tabular environments are set to be single-spaced in the  thesis class,  but long tables do not use tabular
% % to get around this, set the spacing to single spacing at the start of the long table environment, and set it back to double-spacing at the end of it

% \begin{longtable}{c|c|c}
% \caption[This is the title I want to appear in the List of Tables]{This Is a Table Example} \label{tab:pfams} \\
% \hline
% A & B & C \\
% \hline
% \endfirsthead
% \multicolumn{3}{@{}l}{} \\
% \hline
% A & B & C\\
% \hline
% \endhead
% a1 & b1 & c1 \\
% a2 & b2 & c2\\
% a3 & b3 & c3\\
% a4 & b4 & c4\\
% \hline
% \end{longtable}

% The package ``upgreek'' allows us to use non-italicized lower-case greek letters. See for yourself: $\upbeta$, $\bm\upbeta$, $\beta$, $\bm\beta$. Next is a numbered equation:
% \begin{align}
% \label{eq:name}
% \|\bm{X}\|_{2,1}={\underbrace{\sum_{j=1}^nf_j(\bm{X})}_{\text{convex}}}=\sum_{j=1}^n\|\bm{X}_{.,j}\|_2
% \end{align}
% The reference to equation (\ref{eq:name}) is clickable. 
% \section[Theorems, Corollaries, Lemmas, Proofs, Remarks, Definitions and Examples]{Theorems, Corollaries, Lemmas, Proofs, Remarks, Definitions,and Examples}

% \begin{theorem}
% \label{thm:onlytheorem}
% \blindtext
% \end{theorem}

% \begin{proof}
% I'm a (very short) proof.
% \end{proof}

% \begin{lemma}
% I'm a lemma.
% \end{lemma}

% \begin{corollary}
% I include a reference to Thm. \ref{thm:onlytheorem}.
% \end{corollary}

% \begin{proposition}
% I'm a proposition.
% \end{proposition}

% \begin{remark}
% I'm a remark. 
% \end{remark}

% \begin{definition}
% I'm a definition. I'm a definition. I'm a definition. I'm a definition. I'm a definition. I'm a definition. I'm a definition. I'm a definition. I'm a definition. I'm a definition. I'm a definition. 
% \end{definition}

% \begin{example}
% I'm an example.
% \end{example}


% \section[Optional table of contents heading]{Section with\\ linebreaks in\\the
% name}


% \Blindtext[2]




